\documentclass{article}

\usepackage{amsmath,amssymb}

\title{Domain decomposition}
\author{Octave Litrico}
\date{May 2026}

\begin{document}

\maketitle

% ===================== PROBLÈME 1 =====================

\begin{equation}
\begin{cases}
-\Delta u_1^k = f & \text{sur } \Omega_1 \\[4pt]
u_1^k = g & \text{sur } \partial \Omega_1 \setminus (\{\beta\}\times(0,1)) \\[4pt]
u_1^k(\beta,y) = u_2^{k-1}(\beta,y)
\end{cases}
\end{equation}

\begin{equation}
\begin{cases}
-\Delta u_2^k = f & \text{sur } \Omega_2 \\[4pt]
u_2^k = g & \text{sur } \partial \Omega_2 \setminus (\{\alpha\}\times(0,1)) \\[4pt]
u_2^k(\alpha,y) = u_1^k(\alpha,y)
\end{cases}
\end{equation}

% ===================== ERREUR =====================

On pose
\[
e_i^k = u - u_i^k
\]

\begin{equation}
\begin{cases}
-\Delta e_1^k = 0 & \text{sur } \Omega_1 \\[4pt]
e_1^k = 0 & \text{sur } \partial \Omega_1 \setminus (\{\beta\}\times(0,1)) \\[4pt]
e_1^k(\beta,y) = e_2^{k-1}(\beta,y)
\end{cases}
\label{eq:def1}
\end{equation}

\begin{equation}
\begin{cases}
-\Delta e_2^k = 0 & \text{sur } \Omega_2 \\[4pt]
e_2^k = 0 & \text{sur } \partial \Omega_2 \setminus (\{\alpha\}\times(0,1)) \\[4pt]
e_2^k(\alpha,y) = e_1^k(\alpha,y)
\end{cases}
\end{equation}

% ===================== BASE =====================

On pose pour tout $n \in \mathbb{N}$ :
\[
z_n(y) = \sqrt{2}\sin(n\pi y)
\]

La famille $(z_n)_{n\in\mathbb{N}}$ forme une base hilbertienne de $L^2(0,1)$.

% ===================== DÉCOMPOSITION =====================

\[
e_i^k(x,y) = \sum_{n=1}^{\infty} e_{i,n}^k(x)\, z_n(y)
\]



En appliquant \eqref{eq:def1}, on obtient :
\begin{align*}
\sum_{n=1}^{\infty}\Big(\partial_{xx}e_1^k(x) - n^2\pi^2 e_{1,n}^k z_n(y)\Big) &= 0
\end{align*}

Donc
\[
\partial_{xx} e_{1,n}^k - n^2\pi^2 e_{1,n}^k = 0
\]

Donc
\[
e_{1,n}^k(x) = A_n e^{n\pi x} + B_n e^{-n\pi x}
\]

\begin{align*}
\text{or } \quad & e_1(0, y) = 0 \\[6pt]
\text{donc} \quad & \sum_{n=1}^{\infty}(A_n + B_n)\sin(n\pi y) = 0 \\[6pt]
\text{donc} \quad & A_n = -B_n \\[6pt]
\text{en } x=\beta : \quad &
\sum_{n=1}^{\infty} A_n\big(e^{n\pi\beta} - e^{-n\pi\beta}\big)\sin(n\pi y)
= e_2^{k-1}(\beta,y)
\end{align*}

Donc
\[
A_n\big(e^{n\pi\beta} - e^{-n\pi\beta}\big)
= e_{2,n}^{k-1}(\beta)
\]

Donc
\[
A_n =
\frac{e_{2,n}^{k-1}(\beta)}
{e^{n\pi\beta} - e^{-n\pi\beta}}
\]

\vspace{6pt}

De plus, en appliquant le même raisonnement à $e_2^k$, on obtient :
\[
e_{2,n}^k = C_n e^{n\pi x} + D_n e^{-n\pi x}
\]

or
\[
e_2^k(b,y)=0 \quad \forall y\in(0,1)
\quad \text{donc} \quad
D_n = -C_n e^{2n\pi b}
\]

\vspace{6pt}

et $e_2^k(\alpha,y)=e_1^k(\alpha,y)$ donc
\begin{align*}
&\sum_{n=1}^{\infty} C_n\big(e^{n\pi\alpha} - e^{n\pi(2b-\alpha)}\big)\sin(n\pi y) \\
&= \sum_{n=1}^{\infty} A_n\big(e^{n\pi\alpha} - e^{-n\pi\alpha}\big)\sin(n\pi y)
\end{align*}

Donc
\[
C_n =
\frac{e_{2,n}^{k-1}(\beta)}
{e^{n\pi\beta} - e^{-n\pi\beta}}
\cdot
\frac{e^{n\pi\alpha} - e^{-n\pi\alpha}}
{e^{n\pi\alpha} - e^{n\pi(2b-\alpha)}}
\]

Enfin
\[
e_2^k(\beta,y)
=
\sum_{n=1}^{\infty}
e_{2,n}^{k-1}(\beta)
\frac{e^{n\pi\alpha} - e^{-n\pi\alpha}}
{e^{n\pi\beta} - e^{-n\pi\beta}}
\frac{e^{n\pi\beta} - e^{n\pi(2b-\beta)}}
{e^{n\pi\alpha} - e^{n\pi(2b-\alpha)}}
\]
\end{document}